\documentclass[twoside,11pt]{article}

% ── JMLR style (jmlr2e.sty must be in the same directory / Overleaf project).
%    jmlr2e already loads: amssymb, natbib, graphicx, hyperref, epsfig, and
%    defines the theorem/lemma/corollary/definition/proof/example environments.
%    Do NOT re-load those or re-declare \newtheorem (it errors). ──
\usepackage[utf8]{inputenc}
\usepackage[T1]{fontenc}
\usepackage{jmlr2e}

% ── Additional packages (loaded AFTER jmlr2e) ──
\usepackage{amsmath}   % align/equation* etc. (jmlr2e provides amssymb only)
\usepackage{booktabs}
\usepackage{listings}
\usepackage{xcolor}

% ── JMLR running headers (volume/year/pages auto-filled by editor on accept) ──
\usepackage{lastpage}
\jmlrheading{1}{2026}{1-\pageref{LastPage}}{6/26}{}{00-0000}{Nikita Gorshkov}
\ShortHeadings{Alpha-Compositing $\supsetneq$ Softmax Attention}{Gorshkov}
\firstpageno{1}

% ── Lean listing style ──
\lstdefinelanguage{Lean}{
  keywords={theorem,lemma,def,variable,import,open,noncomputable,section,end,by,exact,use,apply,rw,simp,unfold,show,fun,let,have,where,if,then,else,return,do,for,in},
  sensitive=true,
  morecomment=[l]{--},
  morecomment=[s]{/-}{-/},
  morestring=[b]",
}
\lstset{
  language=Lean,
  basicstyle=\ttfamily\footnotesize,
  keywordstyle=\color{blue}\bfseries,
  commentstyle=\color{gray}\itshape,
  breaklines=true,
  frame=single,
  xleftmargin=1em,
  numbers=left,
  numberstyle=\tiny\color{gray},
}

\begin{document}

\title{On the Expressiveness of Alpha-Compositing:\\
A Strict Superset of Softmax Attention}

\author{\name Nikita Gorshkov \email ngorshkov@proton.me \\
       \addr Senefelderstr.\ 29a\\
       10437 Berlin, Germany}

% Action editor is assigned by JMLR; leave the placeholder until then.
\editor{}

\maketitle

% ════════════════════════════════════════════════════════════
\begin{abstract}%
We establish a strict inclusion between two fundamental weighted aggregation schemes used in modern AI and computer graphics. Specifically, we prove that (1) the set of weight vectors achievable by softmax attention is a proper subset of those achievable by alpha-compositing, and (2) provide a constructive mapping from any softmax weight vector to equivalent alpha-compositing parameters via a telescoping tail-sum formula. The converse fails because alpha-compositing can produce exact-zero weights (hard sparsity) and sub-unity weight sums (residual transmittance), both structurally impossible under softmax. Both theorems are machine-verified in Lean~4 with Mathlib, requiring zero unproven assertions and depending only on standard axioms. This result formally grounds the use of volume-rendering equations as drop-in replacements for attention in neural architectures, confirming that no expressiveness is lost and strict additional capability is gained.
\end{abstract}

\begin{keywords}
  attention mechanisms, alpha-compositing, volume rendering, expressiveness, formal verification (Lean~4)
\end{keywords}

% ════════════════════════════════════════════════════════════
\section{Introduction}

\subsection{The Two Schemes}

Modern AI and computer graphics each use a weighted aggregation scheme to combine information from multiple sources into a single output. Despite being developed independently for different purposes, these two schemes turn out to be mathematically related, in a way that, to our knowledge, has not been formally established before.

\paragraph{Scheme A: Softmax Attention (Transformers).}
Every large language model uses the transformer architecture~\cite{vaswani2017}. At its core, a transformer computes outputs as weighted sums:
\begin{equation}
\text{output} = \sum_i w_i \cdot v_i, \quad w_i = \frac{\exp(s_i)}{\sum_j \exp(s_j)}
\end{equation}
Properties of softmax weights: (i) always strictly positive, $w_i > 0$ for all $i$; (ii) always sum to exactly 1, $\sum_i w_i = 1$; (iii) every element contributes at least something.

\paragraph{Scheme B: Alpha-Compositing (Volume Rendering).}
3D Gaussian Splatting~\cite{kerbl2023} and earlier volume rendering techniques~\cite{porter1984,max1995} compute outputs as:
\begin{equation}
\text{output} = \sum_i w_i \cdot v_i, \quad w_i = a_i \cdot T_i, \quad T_i = \prod_{j < i}(1 - a_j)
\end{equation}
Each $a_i \in [0,1]$ is the opacity of element $i$. The transmittance $T_i$ represents how much capacity remains after all preceding elements have contributed. Properties: (i) non-negative, $w_i \geq 0$ (can be exactly zero); (ii) sum to at most 1, $\sum_i w_i = 1 - T_{n+1} \leq 1$; (iii) elements can be completely ignored.

\subsection{The Question}

Both schemes produce a weighted sum of values. Both are differentiable and used in gradient-based optimization. But what is their mathematical relationship? Specifically: (1) Can alpha-compositing reproduce any weights that softmax can? (2) Can softmax reproduce any weights that alpha-compositing can? (3) Or is neither a subset of the other?

\subsection{Related Work}

Linear attention approximations~\cite{katharopoulos2020} have explored connections between attention and kernel methods. Ramsauer et al.~\cite{ramsauer2021} related attention to Hopfield networks. To our knowledge, no prior work has formally proven the strict set-theoretic inclusion between softmax weight vectors and alpha-compositing weight vectors, nor provided the constructive telescoping mapping.

% ════════════════════════════════════════════════════════════
\section{Main Results}

\begin{theorem}[Alpha-compositing $\not\subseteq$ Softmax]
\label{thm:alpha-not-subset}
There exist alpha-compositing weights that no softmax can produce.
\end{theorem}

\begin{proof}
Consider $\mathbf{a} = (0, 1)$. The alpha-compositing weights are $w_1 = 0 \cdot 1 = 0$ and $w_2 = 1 \cdot (1 - 0) = 1$, giving weight vector $(0, 1)$.

But softmax weights are always strictly positive: $w_i = \exp(s_i)/\sum_j\exp(s_j) > 0$ since $\exp(s_i) > 0$ for all $s_i \in \mathbb{R}$. Therefore $(0,1)$ is achievable by alpha-compositing but not by softmax.
\end{proof}

\begin{theorem}[Softmax $\subseteq$ Alpha-compositing, constructive]
\label{thm:softmax-subset}
For any softmax weight vector $\mathbf{w}$ (from any scores $\mathbf{s}$), there exist alpha-compositing parameters $\mathbf{a}$ such that the alpha-compositing weights exactly equal $\mathbf{w}$.
\end{theorem}

\begin{proof}
Given softmax weights $w_1, \ldots, w_n$ (all positive, summing to 1), define:
\begin{equation}
a_i = \frac{w_i}{\sum_{j \geq i} w_j} = \frac{w_i}{R_i}
\end{equation}
where $R_i = \sum_{j \geq i} w_j$ is the tail sum. Note $R_1 = 1$ and $R_i = w_i + R_{i+1}$.

The transmittance through the first $i-1$ elements is:
\begin{equation}
T_i = \prod_{j < i}(1 - a_j) = \prod_{j < i}\frac{R_j - w_j}{R_j} = \prod_{j < i}\frac{R_{j+1}}{R_j}
\end{equation}

This is a telescoping product:
\begin{equation}
T_i = \frac{R_2}{R_1} \cdot \frac{R_3}{R_2} \cdots \frac{R_i}{R_{i-1}} = \frac{R_i}{R_1} = R_i
\end{equation}

Therefore the alpha-compositing weight at position $i$ is:
\begin{equation}
w_i^{\text{alpha}} = a_i \cdot T_i = \frac{w_i}{R_i} \cdot R_i = w_i
\end{equation}
\end{proof}

\begin{corollary}[Strict Inclusion]
\begin{equation}
\{\text{softmax weight vectors}\} \subsetneq \{\text{alpha-compositing weight vectors}\}
\end{equation}
Alpha-compositing can additionally produce weight vectors with exact zeros and weight vectors summing to less than 1 (residual transmittance).
\end{corollary}

% ════════════════════════════════════════════════════════════
\section{Discussion}

\subsection{Implications for AI Architecture Design}

This theorem establishes that the rendering equation from computer graphics is provably at least as expressive as the attention mechanism in transformers for computing weighted aggregations. Any computation a transformer performs through attention weights can be exactly replicated by alpha-compositing, plus alpha-compositing offers additional capabilities (exact sparsity, sub-unity sums) that softmax structurally cannot provide.

This is relevant to architectures that replace transformer attention with the alpha-compositing rendering equation from 3D Gaussian Splatting, as explored in the Semantic Gaussian Splatting (SGS) code repository~\cite{sgs-repo}. The theorem proves this replacement does not lose expressiveness.

\subsection{Implications for Computer Graphics}

The standard alpha-compositing pipeline used in rendering is strictly more powerful than softmax attention as a weighted aggregation scheme. This provides theoretical backing for the empirical success of volume rendering and splatting-based methods over attention-based alternatives in 3D reconstruction tasks.

\subsection{The Extra Expressiveness}

\paragraph{Exact zeros (sparsity).} In softmax, every element always contributes at least $\varepsilon > 0$ to every output. In alpha-compositing, setting $a_i = 0$ means element $i$ contributes exactly zero: true hard sparsity without approximation.

\paragraph{Sub-unity sums.} Softmax weights always sum to exactly 1. Alpha-compositing weights sum to $1 - T_{n+1}$, which can be less than 1 when residual transmittance remains. The residual is a natural uncertainty measure: ``I have only accounted for 80\% of the meaning; 20\% is unresolved.'' Softmax has no such mechanism.

% ════════════════════════════════════════════════════════════
\section{Conclusion}

We have proven that the set of weight vectors achievable by softmax attention is a strict subset of those achievable by alpha-compositing. The proof is constructive: given any softmax weights, the telescoping tail-sum formula $a_i = w_i / \sum_{j \geq i} w_j$ produces alpha-compositing parameters that exactly recover those weights. The result is machine-verified in Lean~4 with zero unproven assertions. This formally grounds architectures that replace attention with volume-rendering equations, confirming no expressiveness is lost.

% ════════════════════════════════════════════════════════════
% JMLR requires a funding + competing-interests disclosure.
\acks{The author received no specific funding for this work and declares no
competing interests.}

% ════════════════════════════════════════════════════════════
\bibliographystyle{plainnat}
\begin{thebibliography}{9}

\bibitem{vaswani2017}
A.~Vaswani, N.~Shazeer, N.~Parmar, J.~Uszkoreit, L.~Jones, A.N.~Gomez, L.~Kaiser, and I.~Polosukhin.
\newblock Attention is all you need.
\newblock In \emph{NeurIPS}, 2017.

\bibitem{kerbl2023}
B.~Kerbl, G.~Kopanas, T.~Leimk{\"u}hler, and G.~Drettakis.
\newblock 3{D} {G}aussian splatting for real-time radiance field rendering.
\newblock \emph{ACM Trans.\ Graph.\ (SIGGRAPH)}, 42(4), 2023.

\bibitem{porter1984}
T.~Porter and T.~Duff.
\newblock Compositing digital images.
\newblock \emph{ACM SIGGRAPH Computer Graphics}, 18(3):253--259, 1984.

\bibitem{max1995}
N.~Max.
\newblock Optical models for direct volume rendering.
\newblock \emph{IEEE TVCG}, 1(2):99--108, 1995.

\bibitem{mildenhall2020}
B.~Mildenhall, P.P.~Srinivasan, M.~Tancik, J.T.~Barron, R.~Ramamoorthi, and R.~Ng.
\newblock {NeRF}: Representing scenes as neural radiance fields for view synthesis.
\newblock In \emph{ECCV}, 2020.

\bibitem{katharopoulos2020}
A.~Katharopoulos, A.~Vyas, N.~Pappas, and F.~Fleuret.
\newblock Transformers are {RNN}s: Fast autoregressive transformers with linear attention.
\newblock In \emph{ICML}, 2020.

\bibitem{ramsauer2021}
H.~Ramsauer, B.~Schafl, J.~Lehner, P.~Seidl, M.~Widrich, et~al.
\newblock Hopfield networks is all you need.
\newblock In \emph{ICLR}, 2021.

\bibitem{demoura2021}
L.~de~Moura and S.~Ullrich.
\newblock The {L}ean 4 theorem prover and programming language.
\newblock In \emph{CADE}, 2021.

\bibitem{sgs-repo}
Semantic {G}aussian {S}platting (SGS).
\newblock Source code repository.
\newblock \url{https://github.com/feamando/sgs}.

\end{thebibliography}

% ════════════════════════════════════════════════════════════
\appendix

\section{Machine-Verified Lean 4 Source}
\label{app:lean}

The following is the complete proof. It compiles with zero \texttt{sorry} (unproven assertions) and depends only on standard axioms (\texttt{propext}, \texttt{Classical.choice}, \texttt{Quot.sound}). Verified by Aristotle (harmonic.fun), Lean~4 v4.28.0, Mathlib v4.28.0.

\begin{lstlisting}
import Mathlib

open Finset BigOperators Real

noncomputable section

variable {n : N}

/-- Softmax attention weight for index `i` given scores `s`. -/
def softmaxWeight (s : Fin n -> R) (i : Fin n) : R :=
  exp (s i) / sum j : Fin n, exp (s j)

/-- Alpha compositing weight for index `i` given opacity values `a`. -/
def alphaWeight (a : Fin n -> R) (i : Fin n) : R :=
  a i * prod j in Finset.Iio i, (1 - a j)

/-- Tail sum: sum_{j >= i} w_j. -/
def tailSum (w : Fin n -> R) (i : Fin n) : R :=
  sum j in Finset.Ici i, w j

/-- Construct alpha compositing parameters from a weight vector. -/
def constructAlpha (w : Fin n -> R) (i : Fin n) : R :=
  w i / tailSum w i

lemma softmaxWeight_pos (s : Fin (n + 1) -> R) (i : Fin (n + 1)) :
    0 < softmaxWeight s i :=
  div_pos (Real.exp_pos _) (Finset.sum_pos
    (fun _ _ => Real.exp_pos _) Finset.univ_nonempty)

lemma softmaxWeight_sum (s : Fin (n + 1) -> R) :
    sum i : Fin (n + 1), softmaxWeight s i = 1 := by
  unfold softmaxWeight
  rw [<- Finset.sum_div _ _ _, div_self <| ne_of_gt <|
    Finset.sum_pos (fun _ _ => Real.exp_pos _) Finset.univ_nonempty]

theorem alpha_not_subset_softmax :
    exists (a : Fin 2 -> R), (forall i, a i in Set.Icc (0 : R) 1) /\
      not (exists (s : Fin 2 -> R), forall i,
        alphaWeight a i = softmaxWeight s i) := by
  use ![0, 1]
  simp_all +decide [Fin.forall_fin_two, alphaWeight]
  exact fun x hx => absurd hx <| ne_of_lt <| softmaxWeight_pos x 0

theorem softmax_subset_alpha (s : Fin (n + 1) -> R) :
    exists (a : Fin (n + 1) -> R), (forall i, a i in Set.Icc (0 : R) 1) /\
      forall i, alphaWeight a i = softmaxWeight s i := by
  apply Exists.intro (fun i => constructAlpha (softmaxWeight s) i)
  apply And.intro
  . exact fun i => constructAlpha_mem_Icc _
      (fun i => softmaxWeight_pos _ _) (by simp [softmaxWeight_sum]) _
  . apply alphaWeight_constructAlpha
    . exact fun i => softmaxWeight_pos s i
    . exact softmaxWeight_sum s

end
\end{lstlisting}

\section{Notation Reference}
\label{app:notation}

\begin{center}
\begin{tabular}{lll}
\toprule
Symbol & Definition & Domain \\
\midrule
$s_i$ & Relevance score for element $i$ & $\mathbb{R}$ \\
$w_i$ (softmax) & $\exp(s_i) / \sum_j \exp(s_j)$ & $(0, 1)$, strictly positive \\
$a_i$ & Opacity of element $i$ & $[0, 1]$, can be zero \\
$T_i$ & Transmittance: $\prod_{j<i} (1 - a_j)$ & $[0, 1]$, non-increasing \\
$w_i$ (compositing) & $a_i \cdot T_i$ & $[0, 1]$, can be zero \\
$R_i$ (tail sum) & $\sum_{j \geq i} w_j$ & $(0, 1]$ \\
\bottomrule
\end{tabular}
\end{center}

\section{Verification Details}
\label{app:verification}

\begin{center}
\begin{tabular}{ll}
\toprule
Property & Value \\
\midrule
Prover & Aristotle (harmonic.fun) \\
Language & Lean 4 v4.28.0 \\
Library & Mathlib v4.28.0 \\
\texttt{sorry} count & 0 \\
Axioms & propext, Classical.choice, Quot.sound (standard) \\
Aristotle Project ID & efb72d79-54a2-49ed-a263-d7b9ce34dc33 \\
Date verified & 2026-04-07 \\
\bottomrule
\end{tabular}
\end{center}

\end{document}
